\documentclass[11pt,a4paper]{article}

% --- Encodage et langue ---
\usepackage[utf8]{inputenc}
\usepackage[T1]{fontenc}
\usepackage[french]{babel}
\usepackage{lmodern}

% --- Mise en page ---
\usepackage[margin=2.5cm]{geometry}
\usepackage{setspace}
\onehalfspacing

% --- Tableaux, figures, maths ---
\usepackage{booktabs}
\usepackage{array}
\usepackage{graphicx}
\usepackage{float}
\usepackage{amsmath,amssymb}
\usepackage{siunitx}
\sisetup{group-separator={\,}, group-minimum-digits=4, detect-weight=true, detect-family=true}

% --- Couleurs et hyperliens ---
\usepackage[dvipsnames]{xcolor}
\usepackage[hidelinks]{hyperref}

% --- Bibliographie simple ---
\usepackage{cite}

% --- En-tête article ---
\title{%
  \textbf{Référence Formantique de l'Orchestre}\\[6pt]
  \large Étude quantitative des convergences spectrales instrumentales\\
  et fondements acoustiques des doublures orchestrales%
}
\author{%
  Yan Maresz\\
  Compositeur, Professeur de nouvelles technologies appliquées à la composition\\
  \textit{Conservatoire National Supérieur de Musique et de Danse de Paris}\\
  \texttt{y.maresz@cnsmdp.fr}%
}
\date{Mars 2026}

\sisetup{detect-weight=true, detect-family=true}
\begin{document}

\maketitle

% ======================================================================
\begin{abstract}
Cette étude présente la première analyse quantitative systématique des
formants spectraux de 30~instruments de l'orchestre symphonique, à partir
de \num{5914}~échantillons de tenues soutenues issus des bases SOL2020
(IRCAM) et d'échantillons d'origines diverses (Yan\_Adds). En introduisant la mesure du \emph{Formant
Principal}~(Fp) --- centroïde spectral pondéré en énergie ---, nous
démontrons que les doublures orchestrales classiques s'expliquent par des
convergences formantiques mesurables. Le résultat central est
l'identification d'un \emph{cluster de convergence} dans la zone
vocalique /o/--/\r{a}/ (377--506~Hz), regroupant 11~instruments de
4~familles différentes avec des écarts aussi faibles que $\Delta =
\SI{0}{Hz}$ (flûte--hautbois) et $\Delta = \SI{11}{Hz}$
(basson--violon, cor--alto). Suite à ces analyses, neuf principes d'orchestration acoustique
sont formulés, validés par un taux de concordance multi-sources de
93\,\% (27/29~instruments comparables). L'ensemble constitue un cadre
quantitatif pour l'analyse et la pratique de l'orchestration.

\medskip
\noindent\textbf{Mots-clés :} formants instrumentaux, orchestration,
analyse spectrale, doublures, centroïde spectral, convergence
formantique, SOL2020, Orchidea.
\end{abstract}


% ======================================================================
\section{Introduction}

Les doublures orchestrales --- l'association de deux instruments jouant
la même ligne mélodique ou des lignes complémentaires --- constituent
l'un des outils fondamentaux de l'écriture orchestrale depuis le
\textsc{xviii}\textsuperscript{e}~siècle. La littérature de traités
d'orchestration (Berlioz, Rimski-Korsakov, Koechlin, Piston, Adler)
prescrit ces associations sur la base de l'expérience auditive
accumulée, sans quantification acoustique systématique.

L'objectif de cette étude est de fournir un \emph{cadre quantitatif}
permettant d'expliquer \emph{pourquoi} certaines doublures fonctionnent
en s'appuyant sur l'analyse des formants spectraux --- les zones de
résonance qui définissent le timbre de chaque instrument. Lorsque deux
instruments sont doublés, leurs spectres se combinent : si leurs
formants coïncident, l'oreille perçoit une couleur timbrale homogène
(\emph{fusion}) ; s'ils divergent, elle perçoit un enrichissement
spectral (\emph{complémentarité}).

L'analyse repose sur \num{5914}~échantillons professionnels couvrant
30~instruments, validée contre quatre sources académiques indépendantes
avec un taux de concordance global de 93\,\%.


% ======================================================================
\section{Corpus et sources}

\subsection{Données primaires}

Le corpus comprend deux sources complémentaires :

\begin{itemize}
  \item \textbf{SOL2020 (IRCAM)} --- 12~instruments solistes,
    \num{2391}~échantillons de tenues soutenues (\emph{sustained}),
    enveloppes spectrales pré-calculées (specenv, 1024~bins).
  \item \textbf{Yan\_Adds} --- 18~instruments supplémentaires (ensembles
    de cordes, bois auxiliaires, cuivres graves), \num{1646}~échantillons
    provenant de VSL, Con~Timbre et enregistrements personnels.
\end{itemize}

Au total, 487~lignes de données SOL2020 et 46~lignes Yan\_Adds (fichiers
CSV~v3) couvrent les 30~instruments analysés.

\subsection{Sources de validation}

Quatre références académiques servent de validation croisée :
Backus~\cite{backus1969} (mesures ciblées cuivres et bois),
Giesler~\cite{giesler1985} (mesures détaillées avec associations
vocaliques explicites, 13~instruments), Meyer~\cite{meyer2009} (analyse
acoustique avec couleurs vocales, zones vocaliques définies), et
McCarty/CCRMA~\cite{mccarty2003} (référence directionnelle, limitations
d'exactitude reconnues par l'auteur). Le taux de concordance global est
de 93\,\% (27 instruments sur 29~comparables).


% ======================================================================
\section{Méthodologie}

\subsection{Périmètre de mesure}

Toutes les valeurs formantiques sont mesurées sur des \textbf{tenues
soutenues} (\emph{sustained}) en jeu \emph{ordinario}, excluant les
transitoires d'attaque, les modes de jeu étendus et les dynamiques
extrêmes. Les techniques incluses sont : normal, arco, sustain,
ordinario, vibrato, non-vibrato, flautando, de \emph{pp} à \emph{ff}.

\subsection{Pipeline d'extraction formantique}

L'extraction des formants repose sur les paramètres suivants : FFT de
\num{4096}~échantillons à \SI{44100}{Hz} (résolution
$\sim$\SI{10.8}{Hz/bin}), plage analysée 100--\SI{3000}{Hz}, détection
de pics sur enveloppe spectrale avec seuil d'amplitude à $-30$~dB sous
le pic maximum et distance minimale de \SI{70}{Hz} entre pics.

Un choix méthodologique critique est la détection du \emph{pic le plus
proéminent} (et non du premier pic), ce qui capture le formant perceptuel
principal et évite les résonances basses secondaires.

\subsection{Note terminologique}

Dans cette étude, F1, F2, \ldots\ F6 désignent les \textbf{pics
spectraux principaux} extraits par analyse d'enveloppe spectrale
(\texttt{scipy.signal.find\_peaks}). Ces pics correspondent aux
résonances principales du système acoustique de l'instrument et sont mis
en regard des voyelles cardinales par \emph{analogie perceptive}, sans
prétendre à une équivalence stricte avec les formants du tractus vocal
humain au sens de Fant~(1960).

\subsection{Le Formant Principal (Fp)}

Pour les instruments à large tessiture, le formant F1 est
\textbf{extrêmement instable} d'un registre à l'autre. Le violon
présente par exemple un F1 dont l'écart-type atteint $\sigma =
\SI{651}{Hz}$ --- soit une variation de plus d'une octave. Cette
instabilité rend F1 peu exploitable comme signature timbrale.

La mesure Fp (\emph{Formant Principal}) résout ce problème. Le Fp est le
centroïde spectral pondéré en énergie dans une bande fréquentielle
optimisée par instrument :

\begin{equation}
  \mathrm{Fp} = \frac{\sum_i f_i \cdot A_i}{\sum_i A_i}
  \label{eq:fp}
\end{equation}

\noindent où $f_i$ est la fréquence du bin~$i$, $A_i$ son amplitude, et
la somme porte sur une bande choisie par instrument (typiquement
600--\SI{1400}{Hz} ou 1000--\SI{2000}{Hz}).

Le gain de stabilité est considérable
(tableau~\ref{tab:fp-stability}) : le Fp est systématiquement plus
stable que F1 ou F2 et correspond au \emph{Hauptformant} décrit par
Meyer~\cite{meyer2009} pour les cordes. Vingt-deux instruments sur
trente bénéficient d'une mesure Fp validée.

\begin{table}[ht]
\centering
\caption{Comparaison de stabilité Fp \emph{vs.}\ F2 (ordinario).}
\label{tab:fp-stability}
\begin{tabular}{lrrrrc}
\toprule
Instrument & Fp (Hz) & $\sigma$ Fp & F2 (Hz) & $\sigma$ F2 & Gain \\
\midrule
Violon         & 893  &  92 & 1\,518 & 651  & $7{,}1\times$ \\
Trompette      & 1\,046 &  98 & 1\,324 & 1\,018 & $10{,}4\times$ \\
Clarinette Sib & 1\,412 & 169 & 1\,130 & 795  & $4{,}7\times$ \\
Cor            & 738  & 112 & 1\,106 & 734  & $6{,}5\times$ \\
\bottomrule
\end{tabular}
\end{table}

L'écart-type $\sigma(\mathrm{Fp})$ reflète indirectement la largeur de
bande de la résonance principale : un $\sigma$ faible indique un formant
étroit et bien défini (forte sélectivité), un $\sigma$ élevé un formant
large et diffus. Ainsi, le violon ($\sigma = \SI{92}{Hz}$) possède un
timbre spectralement plus focalisé que la clarinette~Sib ($\sigma =
\SI{169}{Hz}$).

\subsection{Deux représentations complémentaires par registre}

Pour chaque instrument principal, le document de référence fournit deux
représentations complémentaires, calculées \emph{registre par registre}
(les registres étant définis selon les nomenclatures instrumentales
traditionnelles) :

\begin{enumerate}
  \item Le \textbf{profil formantique moyen} : extraction des formants
    F1--F6 sur chaque échantillon individuel, puis calcul de la médiane
    par registre. Ce profil donne les \emph{positions de référence} de
    chaque pic spectral.
  \item La \textbf{carte spectrale vocalique} : moyenne de toutes les
    enveloppes spectrales du registre, affichée sur un axe fréquentiel
    logarithmique avec superposition des zones vocaliques (/u/, /o/,
    /a/, /e/) et enveloppe cepstrale normalisée. Cette carte montre la
    \emph{forme réelle} du spectre moyen avec ses zones de résonance,
    ses vallées, et le contexte vocalique complet.
\end{enumerate}

Les deux approches convergent pour les formants bien définis (instruments
à anche, cuivres) mais peuvent diverger pour les instruments à grande
tessiture (cordes, flûtes, clarinettes) où la moyenne de l'enveloppe
lisse les pics individuels. Le profil formantique donne les positions de
référence ; la carte spectrale donne le contexte acoustique complet.

Cette double analyse par registre résout le problème de la variabilité
spectrale des instruments à large tessiture : plutôt que de se limiter à
une valeur F1 globale peu informative, l'examen registre par registre
révèle des convergences formantiques qui n'apparaissent que lorsque les
instruments jouent dans un registre spécifique (voir
section~\ref{sec:principes}, principe~8).


% ======================================================================
\section{Correspondance voyelles--fréquences}

Le système vocalique fournit un cadre intuitif pour décrire le timbre
instrumental. D'après Meyer~\cite{meyer2009} et
Giesler~\cite{giesler1985}, les correspondances sont les suivantes
(tableau~\ref{tab:voyelles}).

\begin{table}[ht]
\centering
\caption{Zones vocaliques et instruments représentatifs (F1 ou Fp).}
\label{tab:voyelles}
\begin{tabular}{llll}
\toprule
Voyelle & Fréquence (Hz) & Perception & Instruments \\
\midrule
/u/ & 200--400 & Profondeur & Cb (172), Tuba (226), Cbn (226), Trb (237) \\
/o/ & 400--600 & Plénitude  & Alto (377), Cor (388), Sax alto (398), \\
    &          &            & Cor angl. (452), Cl.~Sib (463), Bn (495), Vln (506) \\
/\r{a}/ & 600--800 & Chaleur & Cl.~Mib (678), Fl (743), Htb (743), Tpt (786) \\
/a/ & 800--1\,250 & Puissance & Cor Fp (738), Vln Fp (893), Tpt Fp (1\,046) \\
/e/ & 1\,250--2\,600 & Clarté & Cl.~Sib Fp (1\,412), Fl Fp (1\,535), Htb Fp (1\,485) \\
/i/ & $>$\,2\,600 & Brillance & Picc. (1\,992 F2), Tpt+harmon (2\,358) \\
\bottomrule
\end{tabular}
\end{table}


% ======================================================================
\section{Résultats par famille instrumentale}

\subsection{Les Bois (plage 226--\SI{2638}{Hz})}

La famille des bois présente la plus grande diversité timbrale de
l'orchestre. Les \textbf{anches doubles} (hautbois, cor anglais, basson,
contrebasson) partagent une zone formantique commune autour de
900--\SI{1200}{Hz} avec une coloration /a/ caractéristique. Les
\textbf{clarinettes} (tuyau cylindrique, harmoniques impairs) n'ont pas
de formant fixe au sens strict : le pic spectral suit la fondamentale ou
le 3\textsuperscript{e}~harmonique. Les \textbf{flûtes} (tuyau ouvert)
présentent un spectre décroissant linéairement au-dessus du fondamental.

Découvertes clés : le cor anglais (F1~=~\SI{452}{Hz}) et le basson
(F1~=~\SI{495}{Hz}) se situent dans la zone /o/, aux côtés du cor
(\SI{388}{Hz}) et du violon (\SI{506}{Hz}). Le cor anglais et la
clarinette~Sib convergent à $\Delta = \SI{11}{Hz}$ (452 \emph{vs.}\
\SI{463}{Hz}). Le basson et le violon convergent également à $\Delta =
\SI{11}{Hz}$ (495 \emph{vs.}\ \SI{506}{Hz}).

\subsection{Les Cuivres (plage 162--\SI{2358}{Hz})}

Les cuivres présentent des formants bien définis et une grande stabilité
spectrale inter-dynamiques (sauf la trompette). Le cor en Fa
(F1~=~\SI{388}{Hz}, zone /o/) est l'instrument cuivre le plus versatile
en termes de doublures, avec un accord unanime des quatre sources
académiques. La trompette en Ut présente un F1 extrêmement variable
($\sigma = \SI{642}{Hz}$), mais son Fp à \SI{1046}{Hz} est $10{,}4$
fois plus stable que F2. Dans le grave, le trombone (\SI{237}{Hz}), le
tuba basse (\SI{226}{Hz}) et le tuba contrebasse (\SI{226}{Hz}) forment
un cluster /u/ très serré ($\Delta \leq \SI{11}{Hz}$), rejoints par le
contrebasson (aussi \SI{226}{Hz}).

\subsection{Les Saxophones}

Le saxophone alto (F1~=~\SI{398}{Hz}) s'inscrit pleinement dans la zone
de convergence /o/--/\r{a}/, aux côtés du cor ($\Delta = \SI{10}{Hz}$)
et de l'alto ($\Delta = \SI{21}{Hz}$). C'est le seul saxophone présent
dans le corpus SOL2020.

\subsection{Les Cordes (plage 172--\SI{1518}{Hz})}

La famille des cordes présente la plus grande variabilité spectrale de
l'orchestre, compensée par la mesure Fp. Le violon
(F1~=~\SI{506}{Hz}, /o/) et le basson (F1~=~\SI{495}{Hz}) convergent à
$\Delta = \SI{11}{Hz}$. Le violoncelle rejoint cette convergence via son
F2~=~\SI{506}{Hz} ($\approx$ F1 basson). L'\textbf{effet de section}
est clairement quantifié : l'ensemble ne déplace pas F1 mais
\emph{lisse} les harmoniques supérieurs (F2 violon solo
$\SI{1518}{Hz} \to$ ensemble \SI{1163}{Hz}, soit $-23\,\%$).


% ======================================================================
\section{Synthèse : le cluster de convergence /o/--/\r{a}/}

Le résultat central de cette étude est l'identification d'un
\textbf{cluster de convergence} dans la zone vocalique /o/--/\r{a}/
(377--\SI{506}{Hz}), regroupant 11~instruments de 4~familles
différentes dans un espace de seulement \SI{129}{Hz}
(tableau~\ref{tab:cluster}).

\begin{table}[ht]
\centering
\caption{Les 6~instruments les plus proches du cluster de convergence.}
\label{tab:cluster}
\begin{tabular}{llrl}
\toprule
Instrument & Famille & F1 (Hz) & $\Delta$ / Violon \\
\midrule
Alto       & Cordes  & 377 & \SI{129}{Hz} \\
Cor        & Cuivres & 388 & \SI{118}{Hz} \\
Sax alto   & Saxophones & 398 & \SI{108}{Hz} \\
Cor anglais & Bois   & 452 & \SI{54}{Hz} \\
Cl.~Sib    & Bois   & 463 & \SI{43}{Hz} \\
Basson     & Bois   & 495 & \SI{11}{Hz} \\
\textbf{Violon} & \textbf{Cordes} & \textbf{506} & --- \\
\bottomrule
\end{tabular}
\end{table}

Les doublures les plus serrées --- basson + violon ($\Delta =
\SI{11}{Hz}$), cor + alto ($\Delta = \SI{11}{Hz}$) --- correspondent
précisément aux associations les plus valorisées de la littérature
orchestrale.

\subsection{Doublures vérifiées}

Le tableau~\ref{tab:doublures} présente les doublures principales
classées par écart $\Delta$ croissant.

\begin{table}[ht]
\centering
\caption{Principales doublures vérifiées (sélection).}
\label{tab:doublures}
\small
\begin{tabular}{llrlrl}
\toprule
Instr.~A & Instr.~B & F1 A & F1 B & $\Delta$ & Qualité \\
\midrule
Flûte        & Hautbois     & 743  & 743  & 0   & Quasi-parfaite $\star$ \\
Tuba CB      & Contrebasson & 226  & 226  & 0   & Quasi-parfaite $\star$ \\
Tuba CB      & Tuba basse   & 226  & 226  & 0   & Quasi-parfaite $\star$ \\
Cor          & Alto         & 388  & 377  & 11  & Quasi-parfaite \\
Cor anglais  & Cl.~Sib      & 452  & 463  & 11  & Quasi-parfaite \\
Violon       & Basson       & 506  & 495  & 11  & Quasi-parfaite \\
Trombone     & Tuba basse   & 237  & 226  & 11  & Quasi-parfaite \\
Trombone     & Contrebasson & 237  & 226  & 11  & Quasi-parfaite \\
Cor anglais  & Cor          & 452  & 388  & 64  & Excellente \\
Trompette    & Hautbois     & 786  & 743  & 43  & Excellente \\
Hautbois     & Basson       & 743  & 495  & 248 & Complémentaire \\
Trompette    & Trombone     & 786  & 237  & 549 & Complémentaire \\
\bottomrule
\end{tabular}
\end{table}

L'analyse par registre révèle en outre des convergences invisibles dans
les données globales. Par exemple, la clarinette~Sib au chalumeau (D3--D4,
F1~=~\SI{215}{Hz}) converge exactement ($\Delta = \SI{0}{Hz}$) avec le
trombone au registre grave ; le cor au médium (F1~=~\SI{323}{Hz})
converge avec l'alto au médium --- confirmant acoustiquement la doublure
emblématique de Brahms.

Il est essentiel de souligner que ces convergences sont valables aussi
bien à l'unisson qu'à l'octave ou à tout autre intervalle. La
convergence formantique étant une propriété du timbre et non de la
hauteur, deux instruments dont les formants coïncident fusionnent
quelle que soit la distance en hauteur entre les notes jouées. C'est
pourquoi les doublures violon--basson ou cor--alto traversent les
registres sans perdre leur efficacité acoustique.


% ======================================================================
\section{Neuf principes d'orchestration acoustique}
\label{sec:principes}

L'analyse des données conduit à formuler neuf principes
fondamentaux.

\paragraph{1. Convergence formantique (fusion).}
Quand deux instruments partagent un F1 similaire ($\Delta \leq
\SI{80}{Hz}$), leurs timbres \emph{fusionnent} : l'oreille perçoit une
couleur homogène plutôt que deux voix distinctes. Seuils mesurés :
$\Delta = 0$ (unisson formantique), $\Delta \leq 30$~Hz (fusion
quasi-parfaite), $\Delta$ 30--80~Hz (fusion efficace).

\textbf{Point capital :} la convergence formantique est une propriété du
\emph{timbre}, pas de la hauteur. Elle ne dépend pas de l'unisson
mélodique. Basson (F1~=~\SI{495}{Hz}) et violon (F1~=~\SI{506}{Hz})
fusionnent à n'importe quel intervalle harmonique --- à l'unisson, à
l'octave, à la tierce --- parce que leurs zones de résonance coïncident
indépendamment de la note jouée. C'est pourquoi ces doublures
fonctionnent sur plusieurs octaves et constituent un pilier de
l'orchestration classique. De même, le tuba basse et le contrebasson
(tous deux à \SI{226}{Hz}) conservent leur unisson formantique /u/ que
l'un joue à l'octave supérieure ou à l'unisson de l'autre.

\paragraph{2. Complémentarité spectrale (enrichissement).}
Quand $\Delta > \SI{200}{Hz}$, les spectres se complètent sans se
couvrir, produisant un enrichissement large bande inaccessible à chaque
instrument seul. Exemples : trompette (\SI{786}{Hz}) + trombone
(\SI{237}{Hz}), $\Delta = \SI{549}{Hz}$ ; violon + violoncelle, $\Delta
= \SI{301}{Hz}$.

\paragraph{3. Effet de section (homogénéisation).}
Le passage du soliste à l'ensemble ne déplace pas F1 mais \emph{lisse}
les harmoniques supérieurs : F2 du violon baisse de 23\,\% en section, F3
de 16\,\%. Le résultat perceptif est un timbre plus \emph{fondu} ---
c'est la différence de texture entre un quatuor et une section
orchestrale.

\paragraph{4. La sourdine comme transposition timbrale.}
Pour les \textbf{cordes} et le \textbf{cor}, la sourdine abaisse
systématiquement F1 de 6 à 40\,\%. Pour les cuivres, l'effet dépend du
type : les sourdines cup et harmon produisent l'effet \emph{opposé} ---
la sourdine harmon de la trompette déplace F1 de la zone /a/ à la zone
/i/ (+200\,\%), le cas le plus spectaculaire du corpus.

\paragraph{5. Familles formantiques transversales.}
Les regroupements formantiques traversent les frontières organologiques :
la \emph{famille /u/} (contrebasse, tuba, contrebasson, trombone),
la \emph{famille /o/--/\r{a}/} (alto, cor, sax alto, cor anglais,
clarinette, basson, violon) et la \emph{famille /a/} (via les Fp)
sont plus pertinentes que les familles instrumentales traditionnelles
pour prédire les fusions timbrales.

\paragraph{6. La zone /o/--/\r{a}/ comme attracteur universel.}
La concentration de 11~instruments dans 129~Hz autour de
377--\SI{506}{Hz} ne peut être un hasard : cette zone correspond au
centre géométrique du spectre perceptif humain (voix parlée, bande de
présence). L'orchestre s'est historiquement construit autour de cette
convergence.

\paragraph{7. Rôle du Fp dans les doublures à large tessiture.}
Pour les instruments à grande tessiture (violon, trompette, clarinette),
le Fp remplace avantageusement F1 ou F2 dans l'analyse des doublures.
Le gain de stabilité atteint un facteur $10{,}4\times$ pour la
trompette.

\paragraph{8. Convergences par registre.}
L'analyse par registre met en lumière des convergences
\emph{invisibles dans les données globales} : $\Delta = 0$~Hz entre
clarinette~Sib (chalumeau) et trombone (grave) ; entre cor (médium)
et alto (médium) ; entre cor (médium) et trompette (médium).

\paragraph{9. Hiérarchie de la fusion.}
Les données permettent d'établir une hiérarchie quantitative :
$\Delta \leq \SI{30}{Hz}$ (fusion quasi-parfaite, $\leq 0{,}3$~Bark),
$\Delta$ 30--\SI{80}{Hz} (fusion efficace, 0,3--0,7~Bark),
$\Delta$ 80--\SI{200}{Hz} (proximité, 0,7--1,5~Bark),
$\Delta \geq \SI{200}{Hz}$ (complémentarité spectrale, $\geq
1{,}5$~Bark).


% ======================================================================
\section{Discussion}

\subsection{Validité du Fp}

Le Fp centroïde spectral s'avère être la mesure la plus pertinente pour
caractériser le timbre global d'un instrument à large tessiture. Il
surpasse dramatiquement F1 et F2 en stabilité, tout en correspondant au
\emph{Hauptformant} décrit par Meyer pour les cordes. L'écart-type
$\sigma(\mathrm{Fp})$ constitue un indicateur de la sélectivité
spectrale (facteur~Q) de la résonance principale.

\subsection{Limites}

Le corpus ne comprend qu'un seul saxophone (alto). Les saxophones
soprano, ténor et baryton constituent une extension future prioritaire.
Les données McCarty/CCRMA présentent des limitations d'exactitude
reconnues par l'auteur et doivent être utilisées avec précaution comme
source de validation. Enfin, l'étude se concentre sur les tenues
soutenues en jeu ordinario : les transitoires d'attaque, les modes de
jeu étendus et les dynamiques extrêmes, qui modifient radicalement le
spectre, ne sont pas couverts.

\subsection{Implications pour la composition}

L'identification quantitative des familles formantiques transversales
offre aux compositeurs un outil nouveau : la possibilité de
\emph{prédire} les fusions et complémentarités timbrales avant même
d'écrire, en choisissant les instruments selon leur zone vocalique
plutôt que selon leur famille organologique. Le seuil de $\Delta \leq
\SI{80}{Hz}$ pour la fusion et de $\Delta \geq \SI{200}{Hz}$ pour
l'enrichissement spectral fournit des critères opérationnels.


% ======================================================================
\section{Conclusion}

Cette étude établit que les doublures orchestrales classiques reposent
sur des convergences formantiques mesurables. Le cluster de convergence
/o/--/\r{a}/ (377--\SI{506}{Hz}), regroupant 11~instruments de
4~familles, constitue le noyau acoustique de l'orchestre symphonique.
L'introduction du Fp (centroïde spectral) comme métrique timbrale,
combinée à l'analyse par registre et à la validation multi-sources
(93\,\% de concordance), fournit un cadre quantitatif sans précédent
pour l'analyse et la pratique de l'orchestration.

Le document complet de référence, incluant les profils formantiques
détaillés, cartes spectrales vocaliques par registre et matrices de
convergence pour les 30~instruments, est disponible en ligne à l'adresse
\url{https://zseramnay.github.io/Orchestration/}.


% ======================================================================
\begin{thebibliography}{9}

\bibitem{backus1969}
  J.~Backus,
  \textit{The Acoustical Foundations of Music},
  W.\,W.~Norton \& Company, New York, 1969.

\bibitem{giesler1985}
  W.~Giesler,
  \textit{Instrumentation: Ein Hand- und Lehrbuch},
  Breitkopf \& Härtel, Wiesbaden, 1985.

\bibitem{meyer2009}
  J.~Meyer,
  \textit{Acoustics and the Performance of Music},
  5\textsuperscript{e}~édition, Springer-Verlag, 2009.

\bibitem{mccarty2003}
  M.~McCarty,
  ``Vowel space chart for orchestral instruments,''
  CCRMA, Stanford University, 2003.

\bibitem{fant1960}
  G.~Fant,
  \textit{Acoustic Theory of Speech Production},
  Mouton, La Haye, 1960.

\end{thebibliography}

\end{document}
